% arara: clean: {extensions: ['log','aux','synctex.gz','out']}
% arara: pdflatex: { shell: yes }
% arara: pdflatex: { shell: yes }
% arara: clean: {extensions: ['log','aux','synctex.gz','out']}
\documentclass[spanish,addpoints,answers,a4paper]{exam}
\usepackage[T1]{fontenc}
\usepackage[spanish]{babel}
\usepackage{libertine}
\usepackage[tmargin=2cm,rmargin=2cm,lmargin=2cm]{geometry}
\usepackage{amsmath,bm,amssymb}
\usepackage[shortlabels]{enumitem}
\usepackage{siunitx}
\sisetup{
	group-four-digits = true,
	group-separator = {,}
}
\usepackage{multicol}

\usepackage{graphicx}
\graphicspath{{../../../images/}}
\pagestyle{plain}

\usepackage{minted}
\setminted[cpp]{tabsize=2}
\usepackage{hyperref}
\renewcommand{\solutiontitle}{\noindent\textbf{Solución}\par\noindent}
\renewcommand\listingscaption{Listado}
\newcommand{\unmarkedfntext}[1]{%
	\begingroup
	\renewcommand\thefootnote{}\footnote{#1}%
	\addtocounter{footnote}{-1}%
	\endgroup
}
\spanishdatedel
\decimalpoint

\def\LOGOUNI{%
	\begin{picture}(0,0)\unitlength=1cm
	\put (-13.5,-2.5) {\includegraphics[width=2.8cm]{logouni}}
	\end{picture}
}

\def\LOGOCTIC{%
	\begin{picture}(0,0)\unitlength=1cm
	\put (-8,-1.2) {\includegraphics[height=1cm]{logocticblack}}
	\end{picture}
}

\newcommand{\mychar}{%
	\begingroup\normalfont
	\includegraphics[height=\fontcharht\font`\B]{Octocat.png}%
	\endgroup
}

\begin{document}

\begin{center}
	\sffamily\bfseries\scshape
	{\Large UNIVERSIDAD NACIONAL DE INGENIERÍA}\LOGOUNI\\
	Centro de Tecnologías de la Información y Comunicaciones\LOGOCTIC\\
\end{center}

\vspace{.8cm}

\begin{center}\sffamily\bfseries\large
	Programación en \texttt{C++} \\
	TEMA $\bm{4}$: Funciones definidas por el usuario: \mintinline{cpp}{void}.
\end{center}

\vspace{.5cm}
\noindent
\makebox[\textwidth]{Nombres y apellidos:\enspace\hrulefill}
\vspace{0.2in}
\makebox[\textwidth]{Nombres y apellidos del instructor:\enspace MSc. Víctor Melchor Espinoza.\hfill}

\section*{Ejercicios de práctica}

\begin{questions}

\question Elabore un programa que pida ingresar dos enteros $m$ y $n$ positivos y llame a dos funciones \mintinline{cpp}{MCD(int a, int b)} y \mintinline{cpp}{MCM(int a, int b)} que retornan el máximo común divisor y el mínimo común múltiplo de $m$ y $n$ respectivamente.

\question Pida ingresar dos vectores en $\mathbb{R}^3$, sume, reste y multiplique, tanto escalarmente como vectorialmente. Imprima todos los resultados y utilice las funciones \mintinline{cpp}{suma()}, \mintinline{cpp}{resta()}, \mintinline{cpp}{productoEscalar()} y \mintinline{cpp}{productoVectorial()}.

\question Implemente una función que devuelva el factorial de un número entero positivo $n$. La función también debe verificar que $n$ sea entero positivo. Considere el
\begin{solution}
	\begin{listing}[H]
		\footnotesize
		\inputminted{cpp}{1.cc}
		\caption{Programa \texttt{1.cc}.}
		\label{lst:2.1}
	\end{listing}
\end{solution}

\end{questions}

\section*{Ejercicios obligatorios}

\begin{questions}

\question Elabore un programa que pida el valor de los coeficientes $a$, $b$ y $c$ de la ecuación cuadrática:
\begin{equation*}
ax^2+bx+c=0,\quad\forall\,\{a, b, c\}\subset\mathbb{R}, a\neq 0.
\end{equation*}
Y a continuación muestre las raíces de dicha ecuación. Para ello deberá calcular el discriminante $\Delta = b^2-4ac$\linebreak e implementar tres funciones:
\begin{minted}{cpp}
realesDiferentes (float a, float b, float Delta);
realesIguales (float a, float b);
complejas (float a, float b, float Delta);
\end{minted}
que mostrarán las raíces según el valor que asuma $\Delta$.

\begin{solution}
	\begin{listing}[H]
		\footnotesize
		\inputminted{cpp}{1.cc}
		\caption{Programa \texttt{1.cc}.}
		\label{lst:1.9}
\end{listing}
\end{solution}

\question Usa la función del ejercicio anterior para determinar el valor de $e$ usando la siguiente serie.
\begin{equation*}
\sum_{n=0}^\infty\frac{1}{n!}=\frac{1}{1}+\frac{1}{1}+\frac{1}{2}+\frac{1}{6}+\frac{1}{24}+\frac{1}{120}+\cdots=e
\end{equation*}

\begin{solution}
	\begin{listing}[H]
		\footnotesize
		\inputminted{cpp}{2.cc}
		\caption{Programa \texttt{2.cc}.}
		\label{lst:1.10}
\end{listing}
\end{solution}

\end{questions}

\begin{flushright}\bfseries
Centro de Tecnologías de la Información y Comunicaciones (CTIC)\\[2mm]
\today\unmarkedfntext{Código disponible en \href{https://github.com/carlosal1015/Cpp-Programming}{\mychar{}}.}
\end{flushright}

\end{document}